\chapter{Validación y Métricas de Calidad}
\label{ch:validation}
% ==============================================================

\section{Estrategia de validación}

Comparación de la línea detectada contra \textbf{transectos GNSS independientes} del IEL:
puntos \textbf{no usados} en la calibración de homografía.

\section{Métricas}

\begin{equation}
  \text{RMSE} = \sqrt{\frac{1}{N} \sum_{i=1}^{N} d_i^2}
  \qquad
  \text{MAE} = \frac{1}{N} \sum_{i=1}^{N} d_i
\end{equation}

donde $d_i$ es la distancia mínima en metros del punto GNSS $i$ a la polilínea detectada.

\section{Criterios de aceptación}

\begin{table}[H]
\centering
\caption{Umbrales de calidad del sistema}
\begin{tabular}{L{4cm} C{3cm} L{5.5cm}}
\toprule
\textbf{Métrica} & \textbf{Objetivo} & \textbf{Condición} \\
\midrule
RMSE geométrico    & $< 1{,}5$~m & Condiciones estándar (12:00h, sin interferencias) \\
MAE geométrico     & $< 1{,}0$~m & Objetivo secundario \\
Cobertura temporal & $\geq 90\%$ & Imágenes a 12:00h procesadas sin error \\
$\delta_\text{px}$ alineación & $< 3$~px & Drift residual post-homografía \\
\bottomrule
\end{tabular}
\end{table}

\section{Factores de error a monitorizar}

\begin{itemize}[leftmargin=1.5em]
  \item \textbf{Calidad de imagen}: niebla, reflejo solar, lluvia, espuma excesiva.
  \item \textbf{Nivel de marea}: afecta la posición real de la línea (calibrar en diferentes estados de marea).
  \item \textbf{Deriva de cámara}: movimiento del soporte por viento, mantenimiento u otras causas.
  \item \textbf{Precisión GCPs}: exactitud del levantamiento GNSS del IEL (instrumental + propagación de error).
\end{itemize}

\section{Validación en QGIS}

\begin{enumerate}[leftmargin=1.5em]
  \item Configurar CRS del proyecto en \textbf{EPSG:25830}.
  \item Cargar GCPs del IEL como capa de puntos.
  \item Arrastrar los GeoJSON de \texttt{data/processed/lines/} al proyecto.
  \item Usar ``Distancia a la polilínea más cercana'' para calcular $d_i$.
  \item Comparar con ortofoto PNOA (WMS/WMTS del IGN) para validación visual.
\end{enumerate}

% ==============================================================
